% Formulario de Technologies Optiques - IMT Atlantique
\documentclass[10pt,a4paper,landscape]{article}

% Paquetes
\usepackage[utf8]{inputenc}
\usepackage[spanish]{babel}
\usepackage{amsmath}
\usepackage{amsfonts}
\usepackage{amssymb}
\usepackage{graphicx}
\usepackage{geometry}
\usepackage{multicol}
\usepackage{fancyhdr}
\usepackage{xcolor}
\usepackage{array}
\usepackage{booktabs}
\usepackage{siunitx}

\geometry{a4paper, landscape, margin=1cm}

% Configuración
\pagestyle{fancy}
\fancyhf{}
\rhead{Formulario Technologies Optiques}
\cfoot{\thepage}
\setlength{\columnsep}{1cm}
\setlength{\parindent}{0pt}

% Comandos personalizados
\newcommand{\fbox}[1]{\fcolorbox{black}{yellow!20}{#1}}
\newcommand{\secbox}[1]{\colorbox{blue!20}{\textbf{#1}}}

\begin{document}

\begin{center}
\LARGE{\textbf{FORMULARIO - TECHNOLOGIES OPTIQUES}}\\
\large{IMT Atlantique - 2025/2026}
\end{center}

\begin{multicols}{3}
\small

%=============================================================================
\secbox{CONSTANTES FÍSICAS}
%=============================================================================

\begin{tabular}{ll}
$c$ & $3 \times 10^8$ m/s \\
$h$ & $6.626 \times 10^{-34}$ J·s \\
$q$ & $1.602 \times 10^{-19}$ C \\
$\epsilon_0$ & $8.854 \times 10^{-12}$ F/m \\
$\mu_0$ & $4\pi \times 10^{-7}$ H/m \\
\end{tabular}

\vspace{0.3cm}
\textbf{Energía del fotón:}
\begin{equation*}
\boxed{E = h\nu = \frac{hc}{\lambda} = \frac{1.24}{\lambda[\mu m]} \text{ eV}}
\end{equation*}

%=============================================================================
\secbox{FIBRAS ÓPTICAS}
%=============================================================================

\textbf{Ley de Snell-Descartes:}
\begin{equation*}
\boxed{n_1 \sin\theta_1 = n_2 \sin\theta_2}
\end{equation*}

\textbf{Apertura Numérica (NA):}
\begin{equation*}
\boxed{NA = \sqrt{n_1^2 - n_2^2} = n_1\sqrt{2\Delta}}
\end{equation*}

\textbf{Índice relativo:}
\begin{equation*}
\boxed{\Delta = \frac{n_1 - n_2}{n_1}}
\end{equation*}

\textbf{Vector de onda:}
\begin{equation*}
\boxed{k = n\frac{2\pi}{\lambda} = n\frac{\omega}{c}}
\end{equation*}

\textbf{Dispersión modal (SI):}
\begin{equation*}
\boxed{\Delta\tau = \frac{L \cdot n_1 \cdot \Delta}{c}}
\end{equation*}

\textbf{Producto B·L (Step-Index):}
\begin{equation*}
\boxed{B \cdot L < \frac{c}{n_1 \cdot \Delta}}
\end{equation*}

\textbf{Producto B·L (Graded-Index):}
\begin{equation*}
\boxed{B \cdot L \sim 10^{10} \text{ bit/s·km}}
\end{equation*}

%=============================================================================
\secbox{ECUACIONES DE MAXWELL}
%=============================================================================

\begin{align*}
&\boxed{\nabla \times \vec{E} = -\frac{\partial \vec{B}}{\partial t}}\\
&\boxed{\nabla \times \vec{H} = \vec{J} + \frac{\partial \vec{D}}{\partial t}}\\
&\boxed{\nabla \cdot \vec{D} = \rho_v}\\
&\boxed{\nabla \cdot \vec{B} = 0}
\end{align*}

\textbf{Ecuación de onda:}
\begin{equation*}
\boxed{\nabla \times \nabla \times \vec{E} = -\frac{1}{c^2}\frac{\partial^2 \vec{E}}{\partial t^2} - \mu_0 \frac{\partial^2 \vec{P}}{\partial t^2}}
\end{equation*}

%=============================================================================
\secbox{INTERFERÓMETROS}
%=============================================================================

\textbf{Campo eléctrico:}
\begin{equation*}
\boxed{E = E_0 e^{i(\omega t - kz + \phi_0)}}
\end{equation*}

\textbf{Intensidad:}
\begin{equation*}
\boxed{I \propto E^2}
\end{equation*}

\textbf{Patrón de interferencia:}
\begin{equation*}
\boxed{I \propto \cos^2(k_0 \sin\theta \cdot z)}
\end{equation*}

\textbf{Mach-Zehnder - Salidas:}
\begin{align*}
&\boxed{I_{SA} = I_0 \cos^2\left(\frac{\Delta\phi}{2}\right)}\\
&\boxed{I_{SB} = I_0 \sin^2\left(\frac{\Delta\phi}{2}\right)}
\end{align*}

\textbf{Desfase:}
\begin{equation*}
\boxed{\Delta\phi = \frac{2\pi}{\lambda}\Delta l = \frac{2\pi n e}{\lambda}}
\end{equation*}

\textbf{Fabry-Pérot - Transmisión:}
\begin{equation*}
\boxed{I = I_0 \frac{T^2}{1 + R^2 - 2R\cos\Delta\phi}}
\end{equation*}

\textbf{Condición de resonancia:}
\begin{equation*}
\boxed{\Delta\phi = 2k\pi \Rightarrow \frac{2L}{\lambda} = k}
\end{equation*}

\textbf{Intervalo Espectral Libre:}
\begin{equation*}
\boxed{\Delta\nu_{ISL} = \frac{c}{2nL}}
\end{equation*}

\textbf{Finesse:}
\begin{equation*}
\boxed{\text{Finesse} = \frac{ISL}{\Delta\lambda}}
\end{equation*}

\textbf{Condición de Bragg:}
\begin{equation*}
\boxed{\lambda_{Bragg} = 2n\Lambda}
\end{equation*}

%=============================================================================
\secbox{ACOPLADORES}
%=============================================================================

\textbf{Acoplador direccional (2 guías):}
\begin{align*}
&\boxed{P_1(z) = P_0 \cos^2(\kappa z)}\\
&\boxed{P_2(z) = P_0 \sin^2(\kappa z)}
\end{align*}

\textbf{Acoplador 3dB:}
\begin{equation*}
\boxed{z = \frac{\pi}{4\kappa}}
\end{equation*}

\textbf{MMI - Vector de onda:}
\begin{equation*}
\boxed{k_{z\nu} = k_0 n_1 - \frac{(\nu+1)^2 \pi}{4n_1 W^2}}
\end{equation*}

\textbf{MMI - Longitud de auto-imagen:}
\begin{equation*}
\boxed{L_\pi = \frac{4n_1 W^2}{3\lambda}}
\end{equation*}

\textbf{MMI - Auto-imagen:}
\begin{equation*}
\boxed{z = 2k \cdot 3L_\pi}
\end{equation*}

\textbf{MMI - Imagen espejo:}
\begin{equation*}
\boxed{z = (2k+1) \cdot 3L_\pi}
\end{equation*}

%=============================================================================
\secbox{FOTODIODOS}
%=============================================================================

\textbf{Fotocorriente:}
\begin{equation*}
\boxed{I_{ph} = \mathcal{R} \cdot P_{opt}}
\end{equation*}

\textbf{Responsividad:}
\begin{equation*}
\boxed{\mathcal{R} = \eta \frac{q\lambda}{hc} = \eta \frac{q\lambda}{1.24} \text{ [A/W]}}
\end{equation*}
($\lambda$ en $\mu$m)

\textbf{Energía del gap:}
\begin{equation*}
\boxed{h\nu \geq E_g \Rightarrow \text{absorción}}
\end{equation*}

\textbf{APD - Ganancia:}
\begin{itemize}
\item Si: $V \sim$ 500V, $M \sim$ 100
\item Ge: $V \sim$ 30V, $M \sim$ 10
\end{itemize}

\textbf{Ruido térmico:}
\begin{equation*}
\boxed{P_{noise} = i_n^2 \cdot \Delta f}
\end{equation*}

%=============================================================================
\secbox{MODULADORES}
%=============================================================================

\textbf{Chirp - Modulación directa:}
\begin{equation*}
\boxed{[n] = f(t) \Rightarrow n = f(t) \Rightarrow \lambda = f(t)}
\end{equation*}

\textbf{Mach-Zehnder - Desfase:}
\begin{equation*}
\boxed{\Delta\phi = \frac{\pi V(t)}{V_\pi}}
\end{equation*}

\textbf{Tensión $V_\pi$ (LiNbO$_3$):}
\begin{equation*}
\boxed{V_\pi = \frac{\lambda e}{n_e^3 r_{33} L}}
\end{equation*}
($r_{33} = 30.8 \times 10^{-12}$ m/V)

\textbf{MZ - Potencia de salida:}
\begin{equation*}
\boxed{P_{out} = P_{in} \cos^2\left(\frac{\pi V}{2V_\pi}\right)}
\end{equation*}

\textbf{Electro-absorción:}
\begin{equation*}
\boxed{P_{out} = P_0 e^{-\Gamma \alpha(V) L}}
\end{equation*}

\textbf{Factor de Henry (chirp):}
\begin{equation*}
\boxed{\alpha_H = \frac{d\phi/dt}{d|E|/dt}}
\end{equation*}

\begin{tabular}{ll}
Mod. directa & $\alpha_H = 3$ a $7$ \\
MZ simple & $\alpha_H = -0.7$ \\
MZ push-pull & $\alpha_H = 0$ \\
\end{tabular}

%=============================================================================
\secbox{LÁSERES}
%=============================================================================

\textbf{Condición de resonancia:}
\begin{equation*}
\boxed{2kL = 2m\pi}
\end{equation*}

\textbf{ISL (cavidad FP):}
\begin{equation*}
\boxed{\Delta\nu_{ISL} = \frac{c}{2nL}}
\end{equation*}

\textbf{Tiempo de vida del fotón:}
\begin{equation*}
\boxed{\tau_c = \frac{-2L}{c \ln R}}
\end{equation*}

\textbf{Potencia de salida:}
\begin{equation*}
\boxed{P = \frac{N \cdot h\nu}{\tau_c}}
\end{equation*}

\textbf{Ley de Bernard-Duraffourg:}
\begin{equation*}
\boxed{h\nu < F_c - F_v}
\end{equation*}

\textbf{Ganancia:}
\begin{equation*}
\boxed{G = g(\lambda) \cdot (N - N_0)}
\end{equation*}

\textbf{Ecuaciones de tasa:}
\begin{align*}
&\boxed{\frac{dN}{dt} = \frac{i}{eV} - \frac{N}{\tau_n} - gP(N - N_0)}\\
&\boxed{\frac{dP}{dt} = \beta_s \frac{N}{\tau_n} + gP(N - N_0) - \frac{P}{\tau_c}}
\end{align*}

\textbf{Valores típicos:}
\begin{tabular}{ll}
$\beta_s$ & $\sim 10^{-4}$ \\
$\tau_c$ & $\sim 10^{-12}$ s \\
$\tau_n$ & $\sim 10^{-9}$ s \\
$N_0$ & $\sim 10^{18}$ cm$^{-3}$ \\
$g$ & $\sim 10^{-6}$ cm$^3$/s \\
\end{tabular}

%=============================================================================
\secbox{POLARIZACIÓN}
%=============================================================================

\textbf{Ecuación de la elipse:}
\begin{equation*}
\boxed{\left(\frac{E_x}{E_{0x}}\right)^2 + \left(\frac{E_y}{E_{0y}}\right)^2 - 2\frac{E_x E_y}{E_{0x}E_{0y}}\cos\delta = \sin^2\delta}
\end{equation*}

\textbf{Parámetros de Stokes:}
\begin{align*}
&\boxed{I = E_{0x}^2 + E_{0y}^2}\\
&\boxed{Q = E_{0x}^2 - E_{0y}^2 = I_{0°} - I_{90°}}\\
&\boxed{U = 2E_{0x}E_{0y}\cos\epsilon = I_{45°} - I_{135°}}\\
&\boxed{V = 2E_{0x}E_{0y}\sin\epsilon = I_{CD} - I_{CG}}
\end{align*}

%=============================================================================
\secbox{ESPECTROS}
%=============================================================================

\textbf{Perfil Lorentziano:}
\begin{equation*}
\boxed{I_L(\lambda) = \frac{1}{1 + 4\pi^2 c^2 \frac{(\lambda - \lambda_0)^2}{\Delta\lambda_L^2}}}
\end{equation*}

\textbf{Perfil Gaussiano:}
\begin{equation*}
\boxed{I_G(\lambda) = \exp\left(-\frac{(\lambda - \lambda_0)^2}{2\sigma^2}\right)}
\end{equation*}

\textbf{Ancho espectral Gaussiano:}
\begin{equation*}
\boxed{\Delta\lambda_G = \frac{2\sigma}{\sqrt{2\ln 2}}}
\end{equation*}

\textbf{Ancho espectral y coherencia:}
\begin{equation*}
\boxed{\Delta\lambda \approx \frac{c}{L_c}}
\end{equation*}

%=============================================================================
\secbox{DECIBELIOS}
%=============================================================================

\textbf{Potencia en dBm:}
\begin{equation*}
\boxed{P_{dBm} = 10\log_{10}\left(\frac{P_{mW}}{1\text{ mW}}\right)}
\end{equation*}

\textbf{Atenuación en dB:}
\begin{equation*}
\boxed{A_{dB} = 10\log_{10}\left(\frac{P_{in}}{P_{out}}\right)}
\end{equation*}

\textbf{Ley de Beer-Lambert:}
\begin{equation*}
\boxed{I(z) = I_0 e^{-\alpha z}}
\end{equation*}

\textbf{Atenuación lineal:}
\begin{equation*}
\boxed{\alpha_{dB/km} = \frac{10}{L}\log_{10}\left(\frac{P_{in}}{P_{out}}\right)}
\end{equation*}

%=============================================================================
\secbox{VENTANAS DE TRANSMISIÓN}
%=============================================================================

\begin{tabular}{lll}
\toprule
\textbf{Ventana} & \textbf{$\lambda$} & \textbf{Material} \\
\midrule
1ª & 850 nm & Si, GaAs \\
2ª & 1.3 $\mu$m & InP (III-V) \\
3ª & 1.55 $\mu$m & InP (III-V) \\
\bottomrule
\end{tabular}

\vspace{0.2cm}
\textbf{Características:}
\begin{itemize}
\item 1ª: Fibras multimodo
\item 2ª: Mínima dispersión
\item 3ª: Mínima atenuación
\end{itemize}

%=============================================================================
\secbox{ESPECIFICACIONES TÍPICAS}
%=============================================================================

\textbf{Componentes pasivos:}
\begin{tabular}{ll}
IL & < 0.2 dB \\
PDL & < 0.3 dB \\
WDL & < 0.3 dB (30 nm) \\
ORL & > 40 dB \\
$P_{max}$ & > 27 dBm \\
Temp. op. & -5°C a +65°C \\
Temp. alm. & -40°C a +80°C \\
\end{tabular}

%=============================================================================
\secbox{FUENTES DE LUZ}
%=============================================================================

\textbf{LED:}
\begin{itemize}
\item $P$: 0.1-0.5 mW
\item $\lambda$: 850-1300 nm
\item $\Delta\lambda$: 50-150 nm
\item $B_{max}$: cientos Mbit/s
\end{itemize}

\textbf{Láser multimodo (FP, VCSEL):}
\begin{itemize}
\item $P$: 1-10 mW
\item $\lambda_{FP}$: 1300-1550 nm
\item $\lambda_{VCSEL}$: 850 nm
\item $\Delta\lambda$: 2-10 nm
\item $B_{max}$: hasta 10 Gbit/s
\end{itemize}

\textbf{Láser monomodo (DFB, DBR):}
\begin{itemize}
\item $P$: 1-10 mW
\item $\lambda$: 1300-1550 nm
\item $\Delta\lambda$: 0.1-0.5 pm
\item $B_{max}$: hasta 10 Gbit/s
\end{itemize}

%=============================================================================
\secbox{RELACIONES IMPORTANTES}
%=============================================================================

\textbf{Frecuencia-longitud de onda:}
\begin{equation*}
\boxed{c = \lambda \nu}
\end{equation*}

\textbf{Número de fotones:}
\begin{equation*}
\boxed{N_{photons} = \frac{P \cdot t}{h\nu} = \frac{P \lambda}{hc}}
\end{equation*}

\textbf{Densidad de potencia:}
\begin{equation*}
\boxed{I = \frac{P}{A} \text{ [W/m}^2\text{]}}
\end{equation*}

\textbf{Campo-Intensidad:}
\begin{equation*}
\boxed{I = \frac{1}{2}n\epsilon_0 c |E|^2}
\end{equation*}

%=============================================================================
\secbox{WDM}
%=============================================================================

\textbf{Espaciado ITU:}
\begin{equation*}
\boxed{\Delta\lambda = \lambda_{i+1} - \lambda_i = 0.8 \text{ nm}}
\end{equation*}

\textbf{OSNR (Optical SNR):}
\begin{equation*}
\boxed{OSNR = \frac{P_{signal}}{P_{ASE}}}
\end{equation*}

%=============================================================================
\secbox{RÉCORDS}
%=============================================================================

\begin{itemize}
\item \textbf{Fibra monomodo}: 319 Tbit/s (3000 km, 2021)
\item \textbf{Fibra multimodo}: 1 Pbit/s (23 km, 2020)
\item \textbf{Progreso}: $10^{10}$ × en 40 años (capacidad × distancia)
\end{itemize}

\end{multicols}

\end{document}
